\section{Introduction}\label{sec:introduction}
The relativistic Vlasov--Maxwell system describes a collisionless plasma
whose particles move under their self-consistent electromagnetic field.
For momentum $v\in\R^3$, write
\[
 q(v)=\sqrt{1+|v|^2},\qquad u(v)=\frac{v}{q(v)}.
\]
The particle density is $f=f(t,x,v)\ge0$; the electric and magnetic fields
are $E=E(t,x)$ and $B=B(t,x)$. We use the one-species normalization
\begin{equation}\label{eq:rvm}
\begin{aligned}
 \partial_t f+u(v)\cdot\nabla_x f+
     (E+u(v)\times B)\cdot\nabla_v f&=0,\\
 \partial_t E-\nabla_x\times B&=-j_f,
 &\partial_t B+\nabla_x\times E&=0,\\
 \nabla_x\cdot E&=\rho_f,
 &\nabla_x\cdot B&=0,
\end{aligned}
\end{equation}
where
\[
 \rho_f(t,x)=\int_{\R^3}f(t,x,v)\,dv,
 \qquad j_f(t,x)=\int_{\R^3}u(v)f(t,x,v)\,dv.
\]
There is no background charge.

A \emph{smooth admissible datum} is a triple $(f_0,E_0,B_0)$ satisfying
\begin{equation}\label{eq:data}
\begin{gathered}
 f_0\in C_c^\infty(\R^3_x\times\R^3_v),\qquad f_0\ge0,\\
 E_0,B_0\in C_b^\infty(\R^3;\R^3)\cap L^2(\R^3;\R^3),\\
 \nabla_x\cdot E_0=\rho_{f_0},\qquad \nabla_x\cdot B_0=0.
\end{gathered}
\end{equation}
Here $C_b^\infty$ means that every spatial derivative is bounded.
A \emph{classical solution} on $[0,T)$ is a $C^1$ triple solving
\eqref{eq:rvm} pointwise, with
$E,B\in C([0,T);L^2_x)$, such that the support of $f$ in $(x,v)$
is contained in a compact set on every compact time subinterval.

\Needspace{12\baselineskip}
\begin{theorem}\label{thm:global}
Every smooth admissible datum \eqref{eq:data} generates a unique global
classical solution of \eqref{eq:rvm}. For every finite $T$, this solution
satisfies
\[
 f\in C^\infty([0,T]\times\R^3_x\times\R^3_v),\qquad
 E,B\in C^\infty([0,T]\times\R^3_x;\R^3),
\]
and $f$ has compact phase-space support on $[0,T]$.
The two divergence constraints hold for all time.
\end{theorem}

Theorem~\ref{thm:global} resolves the large-data global classical
regularity problem for this one-species model in three spatial and three momentum
dimensions. The data have no size or symmetry restriction. The field
assumptions allow a nonzero total charge and its Coulomb tail;
they do not require square-integrability of all field derivatives.
The momentum support bound is on each finite time interval and may
grow with that interval.

\subsection{Historical context}
The distinction between weak existence and classical regularity has
shaped the Cauchy problem for Vlasov--Maxwell. Wollman established local
existence and uniqueness for classical solutions~\cite{Wollman1984}.
DiPerna and Lions developed global weak existence for large
data~\cite{DiPernaLions1989}; a direct formulation for the relativistic
one-species system is given by Rein~\cite{Rein2004Weak}.
These weak existence results do not furnish the regularity and
uniqueness asserted in Theorem~\ref{thm:global}.

For classical solutions, Glassey and Strauss identified bounded
momentum support as a continuation condition~\cite{GlasseyStrauss1986}
and proved global existence for initially dilute
plasmas~\cite{GlasseyStrauss1987}.
Global classical existence was also established by Glassey and
Schaeffer in dimensional reductions: one and one-half dimensions,
two and one-half dimensions, and the planar
problem~\cite{GlasseySchaeffer1990OneHalf,GlasseySchaeffer1997TwoHalf,
GlasseySchaeffer1998PlanarI,GlasseySchaeffer1998PlanarII}.
For the spherically symmetric electrostatic reduction, Glassey and
Schaeffer proved global classical existence for the repulsive relativistic
Vlasov--Poisson system~\cite{GlasseySchaeffer1985Symmetric}.
Rein proved persistence of global existence under small perturbations
of suitably decaying global reference solutions, including nearly
symmetric configurations~\cite{Rein1990Generic}.
These results use restrictions on size, dimension, or the reference
configuration.

The continuation theory also led to new methods. Klainerman and
Staffilani gave a Fourier-based proof of the bounded-momentum
criterion~\cite{KlainermanStaffilani2002}, while Bouchut, Golse and
Pallard gave a physical-space proof using a division lemma for wave
kernels~\cite{BouchutGolsePallard2003}. Luk and Strain showed that it
suffices to control the momentum support after projection onto a
fixed two-dimensional plane~\cite[Theorem~1.2]{LukStrain2014Continuation}.
Moment and spatial-density conditions for continuation were developed by
Pallard~\cite{Pallard2005,Pallard2015Refined}, Sospedra--Alfonso and
Illner~\cite{SospedraAlfonsoIllner2010}, and Luk and
Strain~\cite{LukStrain2016Moments}.
We use the Glassey--Strauss criterion as stated in Luk and
Strain~\cite[Theorem~1.1]{LukStrain2014Continuation};
Section~\ref{cont:continuation} verifies its hypotheses for the
normalization and initial field class considered here.

Small-data theory also developed beyond compact momentum support.
Schaeffer allowed a rapidly decaying momentum tail, retaining compact
spatial support and a suitable smallness
condition~\cite{Schaeffer2004}.
Using a vector-field method with modified commutators, Bigorgne obtained
global existence and sharp asymptotic estimates under weighted smallness
assumptions~\cite{Bigorgne2020}.
His result permits nonzero charge and noncompact particle data.
Wang proved a small-data result without compact support using vector-field
and Fourier methods~\cite{Wang2022SmallData}.
Wei and Yang allow arbitrarily large Maxwell fields with suitable
spatial decay while requiring a sufficiently small weighted particle
density~\cite{WeiYang2021}. Bigorgne subsequently proved global
existence and modified scattering with a small distribution function
and a large Maxwell field~\cite{Bigorgne2025SmallDistribution}.
These theorems preserve a particle-smallness hypothesis.
By prescribing asymptotic data and choosing a sufficiently late starting
time, Bigorgne also constructs particular global solutions with large
particle and field norms~\cite{Bigorgne2023ScatteringMap}. This
construction does not prescribe arbitrary Cauchy data.

Wang's July 2026 two-paper treatment establishes global existence for
large cylindrically symmetric data~\cite{Wang2026I,Wang2026II}.
The first paper develops an iterative smoothing construction and the
second supplies pointwise estimates. Their Theorem~1.1 assumes
$H^s$ fields, $s\ge6$, a weighted $H^s$ particle density, and covariance
of the data under rotations about a fixed axis. It allows localized
noncompact particle data and has no smallness requirement.
The theorem here instead concerns arbitrary compatible compactly
supported particle data, without symmetry, and fields in the class
\eqref{eq:data}. Thus the two data classes should not be identified.

The formulas used below have several classical antecedents.
The retarded-field representation goes back to Glassey and Strauss
\cite{GlasseyStrauss1986}. Luk and Strain organize its most singular
terms using electromagnetic components controlled by energy flux
through light cones~\cite[Proposition~2.2 and Section~3]{LukStrain2014Continuation};
Pallard also uses the orthogonality of the acceleration kernel to the
cone's spatial normal~\cite[Lemma~2.1 and Section~2.3]{Pallard2015Refined}.
The change from retarded phase points to initial particle labels appears
in Lazarovici~\cite[Lemma~9.2]{Lazarovici2016}.
The moving-characteristic cone change of variables is used by Pallard
\cite{Pallard2005}; its proof is also given by Luk and Strain
\cite[Proposition~5.2]{LukStrain2014Continuation}.
We derive the required formulas in the normalization~\eqref{eq:rvm}.

\subsection{Proof strategy}
The central estimate concerns the \emph{signed} change of a particle's
momentum. Along a characteristic $(X(t),V_X(t))$, with
$X'=u(V_X)$ and $V_X'=E(t,X)+u(V_X)\times B(t,X)$, it is the impulse
\[
 V_X(t_2)-V_X(t_1)
 =\int_{t_1}^{t_2}
   \bigl(E(t,X(t))+u(V_X(t))\times B(t,X(t))\bigr)\,dt.
\]
We take the norm after integration. Section~\ref{sec:setup} states the
estimate precisely and shows that it places a lower bound on the time
needed to double the largest particle momentum. These lower bounds have
a divergent sum, which rules out unbounded momentum support in finite
time and allows the continuation criterion to apply.

To estimate the impulse, we follow source particles to their intersections
with the backward light cone of the particle receiving the force.
We divide the source contribution into ranges of momentum, distance, and
angles between the source velocity, receiver velocity, and light ray.
Energy flux through the cone and bounds on the time that source particles
spend in these ranges give the direct force estimates of
Sections~\ref{occ:section} and~\ref{sec:direct}.
At receiver energy comparable to $w$, the absolute-force bound has an
extra factor $\sqrt w$ compared with the signed estimate we need.
The loss occurs when the source velocity is much closer to the light-ray
direction than the receiver velocity is.

In this narrow angular region, we integrate the retarded force along a
source trajectory before taking absolute values.
The identity in Section~\ref{can:section} combines the part without
source acceleration with the part containing it.
The former cancels a term generated by differentiating the cone geometry.
The force is thereby written as a total derivative plus less singular
bulk terms, one of which contains the receiver's acceleration. Expressing
that acceleration through $V_X'$ turns this into an error proportional
to the receiver force. The proof must control its coefficient as well
as the other integration-by-parts errors.

We use the signed identity twice, under an assumed bound on all signed
momentum increments. First, after projection perpendicular to a particle's
momentum, it bounds the number of direction changes of a prescribed
angular size.
This step uses only the preliminary occupation and direct force estimates.
When both source and receiver momenta are large, direction control gives
a stronger occupation bound at intermediate relative angles: between
changes of direction, relative motion makes a monotone crossing of the
region contributing to a fixed range.
Section~\ref{ang:section} proves the direction bound before using it to
improve occupation.

For the second use, Section~\ref{sel:section} selects only those ranges
whose direct force estimates exceed a summable threshold. The other
ranges already satisfy the required bound. On the selected ranges, the
coefficient multiplying the receiver force is at most a constant times
$w^{-1/2}$ after summation. This exactly compensates for the $\sqrt w$
loss in the available absolute-force bound. The occupation improvement
need only hold in the stated angular window: the selection inequalities
force any possible obstruction to this coefficient bound into that
window, where improved occupation rules it out. This yields an impulse
estimate with a time weight that vanishes at the interval endpoints.
Section~\ref{sec:closure}
restores the short endpoint intervals using the assumed signed bound,
obtains a strict improvement, and closes the bootstrap by continuity.

\subsection{Organization}
Section~\ref{sec:setup} states the signed estimate and deduces
Theorem~\ref{thm:global} from it and a continuation result.
Sections~\ref{ret:section}--\ref{sec:direct} derive the retarded force,
occupation estimates and direct bounds.
Section~\ref{can:section} proves the signed cancellation;
Section~\ref{ang:section} bounds direction changes and improves occupation.
Section~\ref{sel:section} estimates the selected signed force, and
Section~\ref{sec:closure} closes the bootstrap.
Section~\ref{sec:continuation-proof} proves local existence, uniqueness,
smooth persistence and the field modification needed to apply the bounded-momentum
continuation theorem to \eqref{eq:data}.
